\chapter{\MakeUppercase{Considerações Finais}}\label{sec:conclusao}

Este capítulo apresenta uma avaliação do produto desenvolvido durante o estágio e uma reflexão sobre o impacto da experiência na formação acadêmica e profissional do autor. Retomam-se, de forma sintética, os objetivos definidos na introdução, discutindo em que medida foram alcançados e quais limitações e aprendizados foram observados ao longo do processo.

\section{Avaliação do Produto Desenvolvido}

O objetivo central do projeto era substituir o registro descentralizado de dados de pesquisa, antes mantido em planilhas eletrônicas, por uma plataforma web própria capaz de centralizar os protocolos da Terra Gerais, controlar o acesso por perfil de usuário e gerar análises visuais de forma automatizada. Diante das funcionalidades entregues, considera-se que esse objetivo foi atendido no plano técnico.

A camada de autenticação e o controle de acesso por perfil garantiram que cada colaborador interagisse apenas com as funcionalidades pertinentes ao seu papel, com renovação transparente de sessão por meio do mecanismo de \textit{refresh token}. O módulo de registros agrícolas implementou o cadastro completo dos protocolos, incluindo formulários extensos com validação declarativa e estruturas dinâmicas para as avaliações de campo. Por fim, o módulo de análise concretizou o principal diferencial da plataforma em relação às planilhas: a geração automatizada de gráficos e \textit{dashboards}, com exportação das visualizações em imagem padronizada pela identidade visual da empresa.

Entre as principais dificuldades superadas, destacam-se a tradução das regras de negócio do domínio agronômico para a lógica dos formulários, que exigiu validação contínua junto ao cliente, e a manutenção de um desempenho adequado na renderização simultânea de múltiplos gráficos, mitigada pela adoção de carregamento sob demanda. Como limitação, observou-se que a adoção da plataforma pelo cliente foi inicialmente reduzida, o que pode ser atribuído à mudança de cultura exigida pela transição das planilhas para um sistema estruturado, e não a uma deficiência técnica da solução. Trabalhos futuros poderiam incluir o acompanhamento dessa adoção ao longo do tempo, a realização de testes de usabilidade com os colaboradores e a evolução dos \textit{dashboards} com novos tipos de análise.

\section{Impacto do Estágio na Formação Acadêmica e Profissional}

O estágio representou a primeira experiência profissional do autor com o desenvolvimento de \textit{software} em um produto real, em equipe estruturada e com um cliente externo. Essa vivência permitiu aplicar, em um contexto concreto, conhecimentos adquiridos ao longo da graduação em disciplinas de engenharia de \textit{software}, desenvolvimento web, banco de dados e interação humano-computador, evidenciando como esses conteúdos se articulam na prática.

No plano técnico, a experiência consolidou competências em desenvolvimento \textit{front-end} com React e TypeScript, no uso de bibliotecas de interface e de gerenciamento de estado, na integração com APIs e, sobretudo, na visualização de dados por meio de gráficos e \textit{dashboards}, área que passou a orientar os interesses profissionais do autor. A adoção da metodologia Scrum proporcionou contato com o trabalho iterativo e incremental, com reuniões de planejamento e acompanhamento, e com a importância das entregas frequentes ao cliente.

No plano interpessoal, o estágio desenvolveu habilidades de comunicação e de trabalho em equipe, especialmente no levantamento de requisitos e na tradução de necessidades de negócio em soluções técnicas, atividade que exigiu diálogo constante com o cliente e com os demais membros da equipe. De modo geral, a experiência contribuiu de forma significativa para o amadurecimento profissional do autor e para a definição de um direcionamento de carreira voltado ao desenvolvimento de interfaces orientadas a dados, integrando a construção de aplicações web à análise e visualização de informações.
